\subsection{Three connecting arrows and constant minimal actions}

\begin{lemma}[Proved: a proper biset retract of a projective plane]
\label{CP2:A-biset-retract}
Let $T\subseteq S$ be nonempty finite left-$N$, right-$M$ bisets,
and let $r:S\to T$ be an equivariant retraction. Write $X_S,X_T$
for their two-object collage spaces, with the same corner monoids.
If $X_S\simeq\CP^2$, then $X_T\simeq\CP^2$. If $T\ne S$,
the latter category has strictly fewer arrows.
\end{lemma}
\begin{proof}
The inclusion and retraction are actual functors, identity on both
corners, and give a retraction $X_S\to X_T$. Use the two-sided
action-bar double mapping cylinders as CW models. Their vertex
spaces are identical and the edge inclusion $E_T\subseteq E_S$
is a simplicial subcomplex. Collapsing $X_T$ therefore gives
$X_S/X_T=\Sigma(E_S/E_T)$, with the evident pointed quotient.
All its positive-degree cohomology products vanish.

If $H^2(X_T;\mathbb Q)=0$, the generator of $H^2(X_S;\mathbb Q)$
lifts from this suspension quotient and has square zero, a
contradiction. Retraction makes every integral homology group
of $X_T$ a direct summand of the corresponding group of $X_S$.
Thus $H_2(X_T;\mathbb Z)=\mathbb Z$, and retraction induces
an integral isomorphism in degree two. Its generator has
primitive nonzero square, since its pullback is a generator
on $X_S$. Consequently $H_4(X_T;\mathbb Z)=\mathbb Z$ and
the degree-four pullback is an isomorphism as well. All other
positive homology groups vanish. The fundamental group of
$X_T$ is a retract of the trivial group, hence trivial.
The projective-plane recognition theorem used in the upper
construction now gives $X_T\simeq\CP^2$. Only the cross-hom
set changed, so the arrow difference is $|S|-|T|$.
\end{proof}

\begin{lemma}[Proved: the complete three-state image classification]
\label{CP2:A-three-state-images}
Two commuting transformation monoids on a three-element set
satisfy at least one of these alternatives: a common fixed
point; both images consist of permutations; one image is
trivial; a proper idempotent transformation commutes with
both images; or, after simultaneous state relabelling and
possibly exchanging the images, they are
\[
 \{1,t_0,t_1\},\qquad \{1,c_0,c_1\},
\]
where $t_i$ fixes $0,1$ and sends $2$ to $i$, while $c_i$
is constant at $i$.
\end{lemma}
\begin{proof}
We provide a fresh complete finite census and a separate
enumeration. Encode a map by $f(0)+3f(1)+9f(2)$, so every
one of the $27$ maps occurs. The producer
\path{code/three_state_images.py} starts with the identity,
adjoins each absent map in turn, and takes multiplication
closure. Every submonoid is reached by successively adjoining
its elements; equality of the complete bit sets removes only
duplicates. This gives exactly $699$ submonoids.

The independent C++ checker
\path{code/check_three_state_images.cpp} uses included and
excluded bit sets instead. It propagates the Horn implications
$a,b\in P\Rightarrow ab\in P$ and their two contrapositives
when the product is excluded, then branches on the least
undecided map. These are necessary closure conditions, and
both branches are tried, so induction on the decision tree
proves coverage. Its $1463$ nodes give exactly the same literal
set of $699$ submonoids; it does not import the producer.

Testing commutation for every ordered pair leaves $1854$
pairs. In the stated order of alternatives the disjoint counts
are $595,5,1158,90,6$. For the fourth alternative every map
is tested for idempotence, proper image and commutation with
every member of both monoids. The checker independently
constructs all displayed residual pairs by choosing the two
core states and exchanging the sides. Their literal set is
exactly the six remaining pairs. Full maps, monoid masks,
pair classifications, witnesses and independent receipts are
in \path{outputs/three_state_images/}. This classifies images
only; no original corner is replaced by its image.
\end{proof}

\begin{lemma}[Proved: a permutation biset has zero second products]
\label{CP2:A-permutation-biset}
If both actions on a nonempty finite biset are permutations,
all rational degree-two products of its collage vanish.
\end{lemma}
\begin{proof}
The two-sided edge bar $E$ projects to $BM\times BN$ as the
finite covering associated to the commuting permutation
actions. One may construct it from the action category of
the product monoid: every transition is a permutation, so
each simplex has exactly $|S|$ compatible lifts, and the
usual face and degeneracy identifications give a covering.
The bisimplicial two-sided bar and this diagonal action bar
have naturally equivalent realizations over $BM\times BN$.
Summing the lifts defines the chain transfer, whose composite
with pullback in rational cohomology is multiplication by
$|S|$. Thus pullback from the product is injective.

Finite monoids have zero rational first cohomology. Kunneth
therefore identifies $H^2(BM\times BN;\mathbb Q)$ with the
direct sum of the two corner second-cohomology groups.
The Mayer--Vietoris compatibility of the two restrictions
of a collage class says that their difference pulls back
to zero on $E$. Injectivity forces both restrictions to
be zero. The relative-cup lemma proves the assertion.
\end{proof}

\begin{theorem}[Proved: compression when the minimal ideal acts constantly]
\label{CP2:A-constant-act-compression}
Let $K$ be the proper minimal ideal of a finite monoid $M$,
and let $I$ be its set of minimal right ideals. Suppose
every element of $K$ acts constantly on a nonempty finite
right $M$-set $S$. There is a finite monoid $N$, with
minimal ideal the rectangular band $I\times S$, such that
\[
 |N|=|M|-|K|+|I||S|,
 \qquad BN\simeq B\mathcal C(M,S).
\]
Here the act cone has source corner $M$, trivial sink
corner, and cross-hom $S$. No triviality assumption is
made on the maximal group of $K$.
\end{theorem}
\begin{proof}
Write $c_k$ for the constant value of the actor $k$.
The right action law gives $c_{k\ell}=c_\ell$ for
$k,\ell\in K$. The minimal-ideal coordinates show that
$c_k$ depends only on its column: left multiplication
by $K$ traverses all row and group coordinates in a
fixed column. Let $\rho(k)$ be its row, and let
$\sigma_m:I\to I$ be the row map induced by left
multiplication by $m$. Then
\[
 \eta(k)=(\rho(k),c_k),\quad
 \eta(mk)=(\sigma_m\rho(k),c_k),\quad
 \eta(km)=(\rho(k),c_km).
\]
Define $N=(M\setminus K)\sqcup(I\times S)$. For old
elements retain their product when it is outside $K$
and otherwise take its $\eta$ image; set
\[
 m(i,s)=(\sigma_m(i),s),\qquad
 (i,s)m=(i,sm),\qquad
 (i,s)(j,t)=(i,t).
\]
The two displayed equivariance equations and the original
associativity check every product involving the new band;
products of three old elements reduce to the corresponding
old associative product and its image. The old identity
is the identity. Thus $q:M\to N$, equal to $\eta$ on
$K$ and the identity on its complement, is a unital
homomorphism. The new band is an ideal and is simple,
and every nonempty ideal meets it, so it is the minimal
ideal of $N$.

Fix $i_0,s_0$ and $e=(i_0,s_0)$. In the full subcategory
$D$ of the Karoubi envelope on $1,e$, the four hom-sets are
\[
 D(1,1)=N,\quad D(e,e)=\{e\},\quad
 D(1,e)=\{i_0\}\times S,\quad
 D(e,1)=I\times\{s_0\}.
\]
Composition is multiplication in $N$, in chronological
categorical order. The idempotent object is a retract
of $1$, hence $BD\simeq BN$. The functor
$F:\mathcal C(M,S)\to D$ sends the two objects to
$1,e$, the source corner by $q$, and $s$ to $(i_0,s)$.
The action equations verify that it preserves composition.
The comma category $(F\downarrow e)$ has terminal
object given by the sink and its identity.

We describe the other comma category, not merely its
homology. Its objects over the source are $n\in N$;
its objects over the sink are $*_i$, $i\in I$.
An arrow $n\to n'$ is labelled $m\in M$ with
$n'q(m)=n$. There is one arrow $(i,s)\to *_i$
labelled $s$, and no other nonidentity arrow involving
a star as source. For a right act $T$, write $A(T)$
for the category whose arrows $x\to y$ labelled $m$
satisfy $ym=x$. Its nerve in degree $d$ is naturally
$T\times M^d$, and so preserves set pushouts and
inclusions. The right-act pushout
$N=M\sqcup_K(I\times S)$ shows that the realization
of our comma category is
\[
 BA(M)\ \sqcup_{BA(K)}\
       \coprod_{i\in I}\operatorname{Cone}(BA(S)).
\]
Indeed each row is a sieve, so a chain ending at $*_i$
has every previous object in that row. There are no
additional mixed simplices. The inclusion $BA(K)\to
BA(M)$ is a simplicial cofibration. Each component
of $BA(K)$ is $BA(e_iM)$ for a minimal idempotent $e_i$.
It is contractible: the arrows $x\to e_i$ labelled
$x$ give a natural transformation to the constant
functor; its naturality is the equation $ym=x$.
The row projection $BA(K)\to I$ is therefore a
homotopy equivalence. Also $A(M)$ has terminal object
$1$. Contracting the cones identifies the displayed
homotopy pushout with a contractible space. Both comma
categories are contractible; Quillen's Theorem A gives
$BF$ a homotopy equivalence and proves the claim.
\end{proof}

\begin{lemma}[Proved: constant-action quotient detects second cohomology]
\label{CP2:A-constant-quotient-H2}
For a finite monoid $M$ with minimal ideal $K$ and nonempty
finite right act $S$, impose the equivalence relation generated
by $sk\sim s'k$ for $s,s'\in S$, $k\in K$. The quotient
$\bar S$ is a right act on which every $k$ is constant, and
the induced act-cone functor gives a surjection
\[
 H^2(B\mathcal C(M,\bar S);\mathbb Q)
 \longrightarrow H^2(B\mathcal C(M,S);\mathbb Q).
\]
If $|S|\leq3$ and some $K$ actor is nonconstant, every
rational degree-two product on the original cone is zero.
\end{lemma}
\begin{proof}
The relation is invariant since $km\in K$. Put
$R=\mathbb QM$, $L=\ker(\mathbb QS\to\mathbb Q)$.
Its relation subspace is exactly $LK$, because it is
spanned by $([s]-[s'])k$; differences along a chain of
generating identifications add, so taking equivalence
closure adds no further linear relation. Thus
$\bar L=\ker(\mathbb Q\bar S\to\mathbb Q)=L/LK$.

Average the maximal group of a minimal idempotent to
obtain the projective left row module $P=Rf\cong\mathbb QI$.
Its augmentation sequence is
$0\to V\to P\to\mathbb Q\to0$, with $KV=0$.
Write $D$ for rational duality. Every $R$-linear map
$V\to DL$ annihilates $LK$: for $v\in V$, $z\in L$,
$k\in K$, its value on $zk$ equals its value at $kv$
on $z$, hence is zero. Consequently
\[
 \operatorname{Hom}_R(V,D\bar L)
   =\operatorname{Hom}_R(V,DL).
\]
The reduced act-cone bar is $(L\otimes_R^{\mathbf L}\mathbb Q)[1]$.
Duality and the projective augmentation sequence identify
its second cohomology with the cokernel of
$\operatorname{Hom}_R(P,DL)\to\operatorname{Hom}_R(V,DL)$.
The quotient act has the same numerator and possibly a
smaller denominator. Naturality proves the stated surjection.
If a $K$ actor is nonconstant, at least two distinct states
are identified, so $|\bar S|\leq2$ when $|S|\leq3$.
Every action on two states extends affinely to the interval,
and its cone products vanish by Lemma~\ref{CP2:A-high-affine}.
The same is true for a singleton. Surjectivity and cup
naturality transfer this conclusion to the original cone.
\end{proof}

\begin{theorem}[Proved: the exact three-state act-cone minimum]
\label{CP2:A-three-state-minimum22}
An ordinary three-state act cone of type $\CP^2$ has at
least twenty-two arrows. The certified twenty-two-arrow
models attain this bound.
\end{theorem}
\begin{proof}
Suppose its monoid has order $m\leq17$. Nonzero degree-two
square requires $H^2(BM;\mathbb Q)\ne0$: otherwise all
second classes lift from the suspension quotient of $BM$.
Writing $|K|=rsh$ for the minimal row, column and group
counts, the rational minimal-group collapse and balanced
tensor formula force $r,s\geq2$. In particular $K$ is
proper, so $q=m-rsh\geq1$.
Lemma~\ref{CP2:A-constant-quotient-H2} forces constant
$K$ action. The compression theorem gives a homotopy-equivalent
monoid with order $m-rsh+3r$ and minimal band $r\times3$.
If $sh\geq3$, its order is at most $m<23$, impossible.
Thus $s=2,h=1$, and the compressed order is $m+r$, where
$r\leq\lfloor(m-1)/2\rfloor$.

For $m\leq15$ this is at most twenty-two. For $m=16$,
order at least twenty-three forces $r=7$, but the compressed
minimal band $7\times3$ differs from both $3\times5$ and
$5\times3$ of the complete minimum-monoid classification.
For $m=17$ the remaining rows are $r=6,7,8$. At $r=6$
the compressed monoid again has order twenty-three with
the wrong minimal-band shape. At $r=7,8$ its contracted
rational algebra $A$ has dimension $q=17-2r=3,1$, and
the row augmentation module has dimension $r-1=6,7$.
The target band--Tor dimensions would be $1,0,1,0,\ldots$.
If $A$ is local, tensor dimension one makes both modules
cyclic by Lemma~\ref{A:local-target-facts}; a cyclic module
cannot have dimension greater than $\dim A$, a contradiction.
If $A$ is nonlocal, Lemma~\ref{A:nonlocal-four} makes nonzero
degree-two Tor persist in arbitrarily high degrees, also
a contradiction. These exhaust all $m\leq17$. Adding
the three states and the sink identity proves the bound.
The already verified upper models have $m=18$ and three
states, proving equality.
\end{proof}

\begin{theorem}[Proved: three connecting arrows reduce to an act cone]
\label{CP2:A-three-cross-reduction}
If a two-object collage with exactly three cross arrows is
of type $\CP^2$, one of its actions is trivial, and its
ordinary three-state act-cone retract is also of type $\CP^2$.
Consequently it has at least twenty-two total arrows,
or at least twenty-three if both corners are nontrivial.
\end{theorem}
\begin{proof}
Apply the complete transformation-image classification.
A common fixed state supplies natural transformations
making both corner inclusions null-homotopic; the
relative-cup lemma excludes this case. Permutation images
are excluded by Lemma~\ref{CP2:A-permutation-biset}.
A proper commuting idempotent gives a biset retract
with at most two states. By Lemma~\ref{CP2:A-biset-retract}
it would still be a projective plane, contrary to the
universal two-cross lemma (or the singleton join).
The residual image pair is excluded by the universal
two-constant theorem, Theorem~\ref{CP2:A-two-constant-image}.

Only a trivial action side remains. Collapsing that
corner is an actual functor to the ordinary act cone,
and adjoining its identity defines a section. By
Lemma~\ref{CP2:A-trivial-biset-side-H2} the induced
map is a rational second-cohomology isomorphism. The
retract therefore retains the nonzero second class.
Its integral homology is a direct summand of that of
the projective plane; in degree two the split injection
is an isomorphism, and its primitive square forces
the same in degree four. Its fundamental group is
trivial by retraction. The recognition theorem gives
the asserted type. Apply the exact three-state act
minimum. Collapsing a corner of order greater than
one saves at least one arrow, proving the last clause.
\end{proof}
